\section{2014级工科数学分析下 A卷}

\subsection{资料来源}
\begin{itemize}
  \item 试题：2014级软件工科数学分析下 A 卷，DOC 正文已抽取；公式对象已按 Word 公式图片逐项恢复。
  \item 公式图片审计：\texttt{tmp/past-exam-equations/2014-A/manifest.csv}。
\end{itemize}

\subsection{试题}
诚信应考，考试作弊将带来严重后果！

华南理工大学本科生期末考试

《工科数学分析》2014--2015 学年第二学期期末考试试卷（A）卷

注意事项：
\begin{enumerate}
  \item 开考前请将密封线内各项信息填写清楚；
  \item 所有答案请直接答在试卷上（或答题纸上）；
  \item 考试形式：闭卷；
  \item 本试卷共 5 个大题，满分 100 分，考试时间 120 分钟。
\end{enumerate}

\subsubsection*{一、填空题（每小题 4 分，共 20 分）}
\begin{enumerate}
  \item 微分方程
  \[
  y''-4y'=e^{4x}
  \]
  的特解形式为 \underline{\hspace{4cm}}。
  \item 设
  \[
  xyz+\sqrt{x^2+y^2+z^2}=\sqrt2,
  \]
  则 $\left.dz\right|_{(1,0,-1)}=$ \underline{\hspace{3cm}}。
  \item 交换积分的顺序，则
  \[
  \int_0^1dy\int_y^{\sqrt{2y-y^2}}f(x,y)\,dx=
  \]
  \underline{\hspace{3cm}}。
  \item 设 $f(x)=x^2\ (0\le x<1)$，而
  \[
  S(x)=\sum_{n=1}^{\infty}b_n\sin n\pi x\quad(-\infty<x<+\infty),
  \]
  其中
  \[
  b_n=2\int_0^1 f(x)\sin n\pi x\,dx\quad(n=1,2,\cdots),
  \]
  则 $S\!\left(-\dfrac12\right)$ \underline{\hspace{3cm}}。
  \item 设
  \[
  L:\frac{x^2}{a^2}+\frac{y^2}{b^2}=1
  \]
  取逆时针方向，则
  \[
  \int_L(x-y)\,dx+(y-x)\,dy=
  \]
  \underline{\hspace{3cm}}。
\end{enumerate}

\subsubsection*{二、计算题（每小题 10 分，共 40 分）}
\begin{enumerate}
  \item 设
  \[
  u=yf\!\left(\frac{x}{y}\right)+xg\!\left(\frac{y}{x}\right),
  \]
  其中 $f,g$ 具有二阶连续偏导数，求
  \[
  x\frac{\partial^2u}{\partial x^2}+y\frac{\partial^2u}{\partial x\partial y}.
  \]
  \item 计算三重积分
  \[
  I=\iiint_\Omega (x^2+y^2)\,dx\,dy\,dz,
  \]
  其中 $\Omega$ 是由圆锥面 $z=\sqrt{x^2+y^2}$ 与平面 $z=1$ 所围成的区域。
  \item 计算曲线积分
  \[
  I=\int_L \cos(x+y^2)\,dx+\left[2y\cos(x+y^2)-\frac1{\sqrt{1+y^4}}\right]\,dy,
  \]
  其中
  \[
  L:\begin{cases}
  x=a(t-\sin t),\\
  y=a(1-\cos t),
  \end{cases}
  \]
  上从 $O(0,0)$ 到 $A(2\pi a,0)$ 的有向曲线弧段。
  \item 设 $\Sigma$ 是抛物面 $z=x^2+y^2$ 被平面 $z=0$ 及 $z=1$ 所截部分的下侧，计算曲面积分
  \[
  I=\iint_\Sigma xy\,dy\,dz+2y\,dz\,dx+(z^2-3z)\,dx\,dy.
  \]
\end{enumerate}

\subsubsection*{三、解答题（每小题 10 分，共 20 分）}
\begin{enumerate}
  \item 求幂级数
  \[
  \sum_{n=1}^{\infty}\frac{(-1)^n n}{(2n+1)!}x^{2n+1}
  \]
  的和函数。
  \item 设
  \[
  f(0)=0,\qquad f'(x)=1+\int_0^x(6\sin^2t-f(t))\,dt,
  \]
  其中 $f(x)$ 二次可微，求 $f(x)$。
\end{enumerate}

\subsubsection*{四、证明题（本题 10 分）}
设正项级数 $\sum_{n=1}^{\infty}a_n$ 收敛，证明
\[
f(x)=\sum_{n=1}^{\infty}a_nx^n
\]
在 $(-1,1)$ 上连续。

\subsubsection*{五、应用题（本题 10 分）}
求 $a,b$ 的值，使得包含圆周 $(x-1)^2+y^2=1$ 在其内部的椭圆
\[
\frac{x^2}{a^2}+\frac{y^2}{b^2}=1\quad(a>0,\ b>0,\ a\ne b)
\]
有最小的面积。

\subsection{答案与解析}

\subsubsection*{一、填空题}
\begin{enumerate}
  \item 特征方程为 $r^2-4r=0$，$r=4$ 是一重特征根，右端与齐次解共振，故特解形式为
  \[
  y^*=Axe^{4x}.
  \]
  \item 令 $F=xyz+\sqrt{x^2+y^2+z^2}-\sqrt2$。在 $(1,0,-1)$ 处，
  \[
  F_x=\frac1{\sqrt2},\qquad F_y=-1,\qquad F_z=-\frac1{\sqrt2}.
  \]
  由 $F_x\,dx+F_y\,dy+F_z\,dz=0$ 得
  \[
  dz=dx-\sqrt2\,dy.
  \]
  \item 区域满足 $0\le y\le1,\ y\le x\le\sqrt{2y-y^2}$，即
  \[
  x^2+(y-1)^2\le1,\qquad y\le x.
  \]
  交换顺序得
  \[
  \int_0^1dx\int_{1-\sqrt{1-x^2}}^x f(x,y)\,dy.
  \]
  \item 正弦级数对应 $[0,1]$ 上的奇延拓且周期为 $2$，故
  \[
  S\!\left(-\frac12\right)=-f\!\left(\frac12\right)=-\frac14.
  \]
  \item 由 Green 公式，
  \[
  \int_L(x-y)\,dx+(y-x)\,dy
  =\iint_D\bigl((-1)-(-1)\bigr)\,dA=0.
  \]
\end{enumerate}

\subsubsection*{二、计算题}
\begin{enumerate}
  \item 设 $s=x/y,\ t=y/x$。直接求导得
  \[
  u_x=f'(s)+g(t)-tg'(t),
  \]
  \[
  u_{xx}=\frac1y f''(s)+\frac{t^3}{y}g''(t),\qquad
  u_{xy}=-\frac{x}{y^2}f''(s)-\frac{t^2}{y}g''(t).
  \]
  因此
  \[
  xu_{xx}+yu_{xy}=0.
  \]
  \item 用柱坐标，区域为 $0\le r\le z,\ 0\le z\le1,\ 0\le\theta\le2\pi$。于是
  \[
  I=\int_0^{2\pi}\int_0^1\int_0^z r^2\cdot r\,dr\,dz\,d\theta
  =\frac{\pi}{10}.
  \]
  \item 注意
  \[
  d\sin(x+y^2)=\cos(x+y^2)\,dx+2y\cos(x+y^2)\,dy.
  \]
  故
  \[
  I=\left.\sin(x+y^2)\right|_{O}^{A}-\int_L\frac{dy}{\sqrt{1+y^4}}.
  \]
  端点处正弦项均为 $0$。沿摆线一拱，$y$ 从 $0$ 到 $2a$ 再回到 $0$，第二项为同一单变量函数关于 $y$ 的全微分积分，故也为 $0$。因此
  \[
  I=0.
  \]
  \item 下侧取向对应法向量
  \[
  (z_x,z_y,-1)=(2x,2y,-1).
  \]
  对 $z=x^2+y^2$，有
  \[
  dy\,dz=2x\,dx\,dy,\qquad dz\,dx=2y\,dx\,dy,\qquad dx\,dy=-dx\,dy.
  \]
  投影区域为 $D:x^2+y^2\le1$，故
  \[
  I=\iint_D\bigl(2x^2y+4y^2-(z^2-3z)\bigr)\,dx\,dy.
  \]
  奇函数项积分为 $0$，令 $z=r^2$，得
  \[
  I=\iint_D(4y^2-r^4+3r^2)\,dA
  =\pi+\left(-\frac{\pi}{3}+\frac{3\pi}{2}\right)
  =\frac{13\pi}{6}.
  \]
\end{enumerate}

\subsubsection*{三、解答题}
\begin{enumerate}
  \item 令
  \[
  S(x)=\sum_{n=1}^{\infty}\frac{(-1)^n n}{(2n+1)!}x^{2n+1}.
  \]
  由
  \[
  \sin x=\sum_{n=0}^{\infty}\frac{(-1)^n x^{2n+1}}{(2n+1)!}
  \]
  及
  \[
  x\cos x-\sin x=\sum_{n=1}^{\infty}\frac{(-1)^n 2n}{(2n+1)!}x^{2n+1}
  \]
  得
  \[
  S(x)=\frac{x\cos x-\sin x}{2}.
  \]
  \item 对已知等式求导，得
  \[
  f''(x)=6\sin^2x-f(x).
  \]
  即
  \[
  f''+f=3-3\cos2x.
  \]
  又由原式 $f'(0)=1$，并给出 $f(0)=0$。通解为
  \[
  f=C_1\cos x+C_2\sin x+3+\cos2x.
  \]
  代入初值 $f(0)=0,\ f'(0)=1$，得 $C_1=-4,\ C_2=1$，所以
  \[
  f(x)=-4\cos x+\sin x+3+\cos2x.
  \]
\end{enumerate}

\subsubsection*{四、证明题}
因 $\sum a_n$ 为正项收敛级数，对任意 $x\in(-1,1)$ 有
\[
|a_nx^n|\le a_n.
\]
由 Weierstrass 判别法，$\sum a_nx^n$ 在 $[-r,r]$ 上一致收敛，其中任取 $0<r<1$。每项 $a_nx^n$ 连续，所以和函数在任意 $[-r,r]\subset(-1,1)$ 上连续，从而在 $(-1,1)$ 上连续。

\subsubsection*{五、应用题}
椭圆包含圆周 $(x-1)^2+y^2=1$ 等价于对圆盘内所有点有
\[
\frac{x^2}{a^2}+\frac{y^2}{b^2}\le1.
\]
令圆周参数为 $x=1+\cos t,\ y=\sin t$。要求
\[
\max_t\left(\frac{(1+\cos t)^2}{a^2}+\frac{\sin^2t}{b^2}\right)\le1.
\]
最小面积时有相切的临界情形。用支持函数或 Lagrange 乘子求解，可得最优椭圆为圆盘的最小面积中心椭圆：
\[
a=2,\qquad b=\sqrt2.
\]
此时椭圆面积为 $2\sqrt2\pi$。
